\section{Công trình liên quan}\label{sec:related}

\paragraph{Quan niệm sai lầm và lỗi logic của người mới học lập trình}
Các tổng quan về lập trình nhập môn mô tả những quan niệm sai về biến, luồng điều khiển, truyền tham số và ``cỗ máy khái niệm'' thực thi chương trình \citep{sorva2013notional,qian2017misconceptions}. Nghiên cứu về lỗi phân biệt sơ suất với lỗi xuất phát từ mô hình sai về ngôn ngữ hoặc đề bài \citep{altadmri2015compilations,hermans2021brain}. \citet{ettles2018logic} phân loại lỗi logic của sinh viên CS1 thành lỗi thuật toán, hiểu sai đề và quan niệm sai, và thấy quan niệm sai là loại phổ biến nhất và khó sửa nhất. Phản hồi dựa trên quan niệm sai có thể giúp sinh viên khi quan niệm sai đã biết trước \citep{gusukuma2018misconception}, còn các công cụ tương tác yêu cầu sinh viên dự đoán hành vi chương trình có thể khơi ra quan niệm sai trực tiếp \citep{lu2024smol}. Các nghiên cứu này thúc đẩy việc gom nhóm lỗi tự động, nhưng cũng cho thấy chỉ riêng chương trình hiếm khi tiết lộ niềm tin nào đã tạo ra nó; vì vậy chúng tôi đánh giá cơ chế mức chương trình và để các khẳng định về nhận thức ra ngoài phạm vi.

\paragraph{Gom cụm chương trình của sinh viên}
OverCode gom lời giải đúng theo mã đã chuẩn hóa \citep{glassman2015overcode}, còn Clara gom lời giải đúng theo tương đương động để sửa bài sai \citep{gulwani2018clara}. Với bài sai, \citet{kaleeswaran2016semisupervised} gom theo chiến lược sửa để giảng viên chỉ sửa một đại diện mỗi cụm; MistakeBrowser gom bài sai theo phép biến đổi mã mà sinh viên đã áp dụng khi tự sửa \citep{head2017mistakebrowser}; TraceDiff đối chiếu vết thực thi của chương trình lỗi với bản sửa gần nhất \citep{suzuki2017tracediff}. Embedding chương trình học từ mã và thực thi đã được dùng để lan truyền phản hồi \citep{piech2015embeddings}, và \citet{shi2021misconception} gom embedding của bài sai dạng khối để giúp chuyên gia phát hiện quan niệm sai. Với C, InvAASTCluster kết hợp bất biến động với cây cú pháp trừu tượng hóa \citep{orvalho2025invaast}, còn MENTOR ghép phân cụm, định vị lỗi bằng công thức và sửa lỗi bằng LLM \citep{orvalho2026mentor}. Các hệ thống này dùng bản sửa hoặc lời giải đúng như một phần của phương pháp. Chúng tôi chỉ dùng bản sửa về sau của chính sinh viên để xây dựng nhãn đánh giá, và hỏi các biểu diễn có sẵn \emph{trước} khi có bản sửa khôi phục nhãn đó tốt đến đâu.

\paragraph{Khai phá quan niệm sai bằng mô hình học}
Các mô hình attention trên cây con định vị lỗi logic chỉ từ nhãn đúng/sai \citep{hoq2025sann}, thành phần tri thức theo mẫu được trích từ mã sinh viên bằng học biểu diễn và phân cụm \citep{edm2026patternkc}, và LLM phát hiện quan niệm sai từ tập mẫu mã trong McMining, với benchmark tiêm quan niệm sai đã biết vào mã Python \citep{alhossami2026mcmining}. Các công cụ lớp học gần đây gom lỗi do LLM nhận diện cho giảng viên \citep{sigcse2026realtime}, định hướng dạy kèm theo các quan niệm sai đã biết \citep{sigcse2026tutor}, hoặc tinh chỉnh tín hiệu phát hiện tự động bằng đánh giá của chuyên gia \citep{aruleba2026ukicer}. PyMETA cho thấy độ chính xác chẩn đoán phụ thuộc mạnh vào cách định nghĩa nhãn chuẩn: mô hình gán tốt lỗi thực thi đầu tiên lại khớp kém hơn nhiều với lộ trình sửa của chuyên gia \citep{li2026pymeta}. Nghiên cứu của chúng tôi cùng mối quan tâm đó và làm cho định nghĩa nhãn trở nên tường minh, kiểm chứng được và tái lập được.

\paragraph{Benchmark, mutant và bản sửa}
C-Pack-IPAs cung cấp bài C90 của 25 bài tập kèm mã sinh viên ẩn danh, test và log chấm lịch sử \citep{orvalho2024cpack}; ITSP là benchmark trước đó gồm chương trình C lỗi và đã sửa \citep{yi2017itsp}. Phân tích đột biến tiêm lỗi có kiểm soát; trong kiểm thử phần mềm, khả năng phát hiện mutant tương quan với khả năng phát hiện lỗi thật \citep{just2014mutants}, còn mutant khái niệm xây từ ví dụ sai của sinh viên làm lộ quan niệm sai về đặc tả \citep{prasad2024conceptual}. Delta debugging cô lập khác biệt tối thiểu gây lỗi giữa một cấu hình lỗi và một cấu hình đạt \citep{zeller2002ddmin}. Chúng tôi kết hợp các ý tưởng này: chạy lại corpus, thu gọn bản sửa của chính sinh viên thành thay đổi tối thiểu đã kiểm chứng, và ghép các nhãn này với mutant một cơ chế có kiểm soát.

\paragraph{Quy nạp luật}
Thuật toán Học Quy nạp (ILA) trích luật sản xuất theo từng lớp từ các ví dụ thuộc tính--giá trị \citep{tolun1998ila}; ILA-2 thêm hệ số phạt để chịu được dữ liệu nhiễu \citep{tolun1999ila2}. ILA đôi khi được nêu cùng các phương pháp phân cụm, nhưng bản chất là quy nạp luật có giám sát, và chúng tôi dùng đúng như vậy: để phát biểu luật NẾU--THÌ ánh xạ bằng chứng thực thi sang cơ chế, và đo xem các luật đó còn đúng tới đâu trên những sinh viên không tham gia quy nạp.

\paragraph{Vị trí của nghiên cứu}
Phần lớn các hệ thống trên phân cụm để tiết kiệm công sức giảng viên hoặc để lan truyền phản hồi, và được đánh giá bằng công sức tiết kiệm được, bằng tính đúng của bản sửa sinh ra, hoặc bằng việc chuyên gia xem xét các cụm. Theo hiểu biết của chúng tôi, chưa công trình nào đo một biểu diễn có sẵn trước mọi bản sửa khôi phục được cơ chế lỗi đã kiểm chứng độc lập tốt đến đâu, và chưa công trình nào báo cáo chữ ký test có thể biểu đạt được bao nhiêu phần cơ chế đó. Đây là hai phép đo mà bài báo này đóng góp, cùng một benchmark để lặp lại được chúng.
